\subsection*{4. The Unit of Analysis Changes the Answer}

Before computing a unit margin, define the \emph{unit}. In education it might be a learner, paid enrollment, completed program, cohort, or institutional contract. In SaaS it might be a customer account, active seat, transaction, or successful workflow. In consulting it may be a project, client, deliverable, or billable expert-hour. In FMCG it might be a pack or an order to a distributor.

A useful unit must connect revenue with the costs and constraints incurred to serve it. A ``₹50,000 contribution per learner'' number is not comparable with a ``₹5 lakh contribution per college contract'' unless we also examine acquisition duration, fixed setup, delivery hours, working capital, and the probability of repeat business. Denominators shape strategic decisions.

Write the unit on top of every table. When units differ, bridge them explicitly: revenue per cohort divided by enrolled learners; contribution per project divided by specialized hours; gross margin per subscription account multiplied by survival probability. Without this discipline, numbers that look comparable are not.

\subsection*{5. The 4 Flows in Every Company}

A business model can be mapped through 4 parallel flows:
\begin{enumerate}
  \item The \textbf{customer flow} runs from discovery through purchase, activation, outcome, and repeat or referral.
  \item The \textbf{work flow} runs from inputs and skills through production, quality control, delivery, and support.
  \item The \textbf{economic flow} runs from revenue to variable costs, fixed overhead, depreciation, financing costs, taxes, and retained profit.
  \item The \textbf{cash flow} runs from deposits, customer collections, supplier payments, payroll, receivables, debt, and investment.
\end{enumerate}

A founder should inspect all 4 because their clocks are different. Marketing may recruit customers in a week; a 6-month program takes half a year to deliver; revenue may be recognized over service time; a college may pay months after examination results. Revenue recognized today may correspond to cash received long ago or cash not yet collected. An operational process can be efficient while its financial cycle is destructive.

\subsection*{6. Revenue Architecture: Who Pays for What, When?}

A company may earn revenue from unit sales, subscriptions, usage charges, transaction commissions, advertisements, licensing, projects, retainers, training cohorts, product-service bundles, or performance-based fees. Each arrangement creates a different relationship between delivered activity and payment.

The 1st danger is confusing buyer payment with beneficiary value. The 2nd is choosing a monetization unit that discourages healthy adoption or hides variable cost. The 3rd is treating recurring billing as recurring customer satisfaction. A large contract is economically impressive only if it is enforceable, collected, fulfilled, and capable of producing a positive contribution under honest labor assumptions.

For each stream record customer, unit, price, discount, recognized revenue logic, collection timing, refunds, service obligations, associated variable costs, acquisition cost, and any renewal or expansion evidence. Where 1 offering supports another, document the cross-subsidy instead of burying it in optimistic consolidation.

\subsection*{7. Gross Merchandise Volume and Platform Turnover Are Not Always Company Revenue}

A marketplace that facilitates ₹10 crore in transactions does not necessarily earn ₹10 crore of revenue. If the platform acts as an agent and retains a 10 percent fee, the relevant recognized revenue may be the fee rather than the gross value; accounting principal-versus-agent assessments depend on specific control and contractual facts under applicable standards. Likewise, a distribution business that buys and resells goods generally has different accounting from an agent arranging 3rd-party sales.

GMV can measure marketplace activity. It is not a synonym for the platform's sales, and the platform's sales are not profit. The founder should distinguish gross flow across an ecosystem from the portion that belongs economically and legally to the company. This is especially important when comparing businesses whose top-line reports use different conventions.
